\documentclass[10pt]{article}
\usepackage[margin=0.7in,top=0.55in,bottom=0.55in]{geometry}
\usepackage{amsmath,amssymb,booktabs,microtype,xcolor,url}
\usepackage[colorlinks=true,linkcolor=blue!60!black,urlcolor=blue!60!black,pdftitle={FRE-GY-9743 Assignment 6},pdfauthor={Raj Pawar}]{hyperref}
\newcommand{\code}[1]{\texttt{#1}}
\newcommand{\E}{\mathbb{E}}
\newcommand{\df}[1]{\mathrm{df}_{#1}}
\newcommand{\usd}{\mathrm{USD}}\newcommand{\eur}{\mathrm{EUR}}\newcommand{\gbp}{\mathrm{GBP}}
\setlength{\parindent}{0pt}\setlength{\parskip}{3pt}\emergencystretch=3em\pagestyle{empty}
\makeatletter\g@addto@macro\normalsize{\setlength{\abovedisplayskip}{5pt}\setlength{\belowdisplayskip}{5pt}\setlength{\abovedisplayshortskip}{3pt}\setlength{\belowdisplayshortskip}{3pt}}\makeatother
\makeatletter\renewcommand{\section}{\@startsection{section}{1}{\z@}{8pt plus 2pt}{3pt}{\normalfont\large\bfseries}}\makeatother
\begin{document}
{\Large\bfseries FRE-GY-9743 Assignment 6: A GBP-Collateralised EUR/USD Forward, and Swap Products}\\[1pt]
Raj Pawar, NYU Tandon School of Engineering, October 28, 2026.\\
Code: \href{https://github.com/pawarraj8888/FRE-GY-9743-Assignments-6/tree/Raj_Pawar_Branch}{github.com/pawarraj8888/FRE-GY-9743-Assignments-6}, branch \code{Raj\_Pawar\_Branch}.\\
Project page: \href{https://pawarraj8888.github.io/FRE-GY-9743-Assignments-6/}{pawarraj8888.github.io/FRE-GY-9743-Assignments-6}.
\vspace{2pt}\hrule

\section*{Part I: an EUR/USD forward collateralised in GBP}
\textbf{Setting.} At $T$ I receive $N$ EUR and pay $NK$ USD, so $K$ is in USD per EUR, like the spot rate $X_t$. The trade
sits under a CSA with cash collateral in GBP, remunerated at SONIA, and I assume perfect collateralisation as in the notes.
I write $c_{\usd}$, $c_{\eur}$, $c_{\gbp}$ for the collateral (OIS) rates SOFR, EuroSTR and SONIA, $r_{\usd}, r_{\eur},
r_{\gbp}$ for the rate at which a desk borrows and invests cash in each currency (one rate per currency, which collapses the
notes' $r_R$ and $r_F$; the gap between those two is the FVA question and I leave it aside), and $y_i := r_i - c_i$ for the gap
between the cash rate and the collateral rate in currency $i$. The $y_i$ are small and not directly observable; they only ever
appear as differences, which is what the cross-currency basis market quotes.

\textbf{Discounting a payoff whose collateral is in another currency.} I redo the replication argument of the collateral notes
(Section 2.2) for a USD payoff $V$ collateralised in GBP. The counterparty posts collateral worth $V$, i.e. $C_{\gbp}=V/X^{\mathrm{GBPUSD}}_t$
pounds, on which I pay SONIA, $c_{\gbp}C_{\gbp}$. I keep the pounds in GBP and invest them at $r_{\gbp}$ (converting them would leave
me short pounds against a liability that is marked in pounds), so the collateral account nets me $(r_{\gbp}-c_{\gbp})C_{\gbp}$ per unit
time, which is $y_{\gbp}V$ in dollars at spot. The derivative itself I fund by borrowing $V$ dollars at $r_{\usd}$, and the hedge is
financed at the secured rate exactly as in Eq.~(12) of the notes. Collecting the carry terms, the self-financing condition becomes
$\partial_tV+\mathcal{L}V = r_{\usd}V - y_{\gbp}V$, and Feynman--Kac gives
\[
V_t=\E^{Q_{\usd}}_t\!\Bigl[e^{-\int_t^T (r_{\usd}-y_{\gbp})\,du}\,V_T\Bigr]
   =\E^{Q_{\usd}}_t\!\Bigl[e^{-\int_t^T (c_{\usd}+y_{\usd}-y_{\gbp})\,du}\,V_T\Bigr].
\]
Two checks: with GBP collateral on a GBP payoff the pounds repay the borrowing directly, the round trip disappears and the rate is
$r_{\gbp}-y_{\gbp}=c_{\gbp}$, SONIA discounting, which is Eq.~(16) of the notes; with USD collateral on a USD payoff it is SOFR. So the general rule is: a payoff in currency $i$ collateralised in
currency $k$ is discounted at $c_i+y_i-y_k$, and I write the corresponding collateralised zero-coupon bond as
\[
P^{(k)}_i(t,T):=\E^{Q_i}_t\!\Bigl[e^{-\int_t^T (c_i+y_i-y_k)\,du}\Bigr],\qquad P^{(i)}_i(t,T)=\E^{Q_i}_t\!\Bigl[e^{-\int_t^T c_i\,du}\Bigr]
\text{ the OIS discount factor.}
\]

\textbf{Q1: valuation formula.} The USD leg is a fixed payment of $NK$ at $T$ under GBP collateral, worth $-NK\,P^{(\gbp)}_{\usd}(t,T)$.
The EUR leg is $N$ euros at $T$ under the same CSA; the argument above run in EUR values it at $N\,P^{(\gbp)}_{\eur}(t,T)$ euros,
and a euro value today is $X_t$ dollars (a spot trade, no funding involved). Hence
\[
V_t = N\Bigl[X_t\,P^{(\gbp)}_{\eur}(t,T)-K\,P^{(\gbp)}_{\usd}(t,T)\Bigr]
    = N\,P^{(\gbp)}_{\usd}(t,T)\bigl[F^{(\gbp)}(t,T)-K\bigr],\qquad
F^{(\gbp)}(t,T):=X_t\,\frac{P^{(\gbp)}_{\eur}(t,T)}{P^{(\gbp)}_{\usd}(t,T)}.
\]
In the language of Section 2.3 of the notes, $P^{(\gbp)}_{\usd}(t,T)$ is the numeraire of the $T$-forward measure for this CSA,
$X_T$ is a martingale under it with $F^{(\gbp)}(t,T)=\E^{T,\gbp}_t[X_T]$, and $V_t=N\,P^{(\gbp)}_{\usd}(t,T)\,\E^{T,\gbp}_t[X_T-K]$.

\textbf{Q2: fair forward strike.} $V_t=0$ gives
\[
K^{*}=F^{(\gbp)}(t,T)=X_t\,\frac{P^{(\gbp)}_{\eur}(t,T)}{P^{(\gbp)}_{\usd}(t,T)},
\]
covered interest parity with both legs discounted on their GBP-collateralised curves. If the spreads $y$ are deterministic the
expectation factorises, $P^{(k)}_i=P^{(i)}_i\,e^{-\int_t^T (y_i-y_k)du}$, and the collateral currency drops out of the ratio:
$K^{*}=X_t\,(P^{(\eur)}_{\eur}/P^{(\usd)}_{\usd})\,e^{-\int_t^T (y_{\eur}-y_{\usd})du}$, the same forward whatever the CSA.
The collateral currency changes how the mark-to-market is discounted, not where the forward is.

\textbf{Q3: are the five discount factors enough?} First what they are. $\df{\mathrm{SONIA}}=P^{(\gbp)}_{\gbp}$,
$\df{\mathrm{EuroSTR}}=P^{(\eur)}_{\eur}$ and $\df{\mathrm{SOFR}}=P^{(\usd)}_{\usd}$ are the OIS curves, cashflows collateralised
in their own currency. $\df{[\mathrm{GBP},\mathrm{USD}]}=P^{(\usd)}_{\gbp}$ and $\df{[\mathrm{EUR},\mathrm{USD}]}=P^{(\usd)}_{\eur}$
are GBP and EUR cashflows collateralised in USD, which is what GBP/USD and EUR/USD cross-currency basis swaps under a USD CSA
bootstrap.\footnote{If the bracket is meant the other way round, USD cashflows collateralised in GBP and in EUR, the same algebra
goes through with the roles swapped: $P^{(\gbp)}_{\usd}=\df{[\mathrm{GBP},\mathrm{USD}]}$ directly and
$P^{(\gbp)}_{\eur}=\df{\mathrm{EuroSTR}}\,\df{[\mathrm{GBP},\mathrm{USD}]}/\df{[\mathrm{EUR},\mathrm{USD}]}$. The conclusion does not change.}
What I need are $P^{(\gbp)}_{\usd}$ and $P^{(\gbp)}_{\eur}$, and neither is on the list, so taken literally the answer is no.
But every $P^{(k)}_i$ differs from the OIS curve $P^{(i)}_i$ only through the collateral switch $e^{-\int (y_i-y_k)}$, and the $y$'s
enter only as differences. Under the assumption that any curve construction already makes when it turns these expectations into
discount factors, that the spreads are deterministic (or at least independent of the rates and of FX), the switch factors chain:
with $Y_{i,k}(T):=e^{-\int_t^T (y_i-y_k)du}$ I have $Y_{i,k}=Y_{i,\usd}\,Y_{\usd,k}$ and $Y_{k,i}=1/Y_{i,k}$, and the two USD bases give
$Y_{\gbp,\usd}=\df{[\mathrm{GBP},\mathrm{USD}]}/\df{\mathrm{SONIA}}$ and $Y_{\eur,\usd}=\df{[\mathrm{EUR},\mathrm{USD}]}/\df{\mathrm{EuroSTR}}$. Then
\[
P^{(\gbp)}_{\usd}=\df{\mathrm{SOFR}}\,Y_{\usd,\gbp}=\df{\mathrm{SOFR}}\,\frac{\df{\mathrm{SONIA}}}{\df{[\mathrm{GBP},\mathrm{USD}]}},\qquad
P^{(\gbp)}_{\eur}=\df{\mathrm{EuroSTR}}\,Y_{\eur,\usd}\,Y_{\usd,\gbp}=\df{[\mathrm{EUR},\mathrm{USD}]}\,\frac{\df{\mathrm{SONIA}}}{\df{[\mathrm{GBP},\mathrm{USD}]}},
\]
and the forward and its value become
\[
K^{*}=X_t\,\frac{\df{[\mathrm{EUR},\mathrm{USD}]}(T)}{\df{\mathrm{SOFR}}(T)},\qquad
V_t=N\,\frac{\df{\mathrm{SONIA}}(T)}{\df{[\mathrm{GBP},\mathrm{USD}]}(T)}\Bigl[X_t\,\df{[\mathrm{EUR},\mathrm{USD}]}(T)-K\,\df{\mathrm{SOFR}}(T)\Bigr].
\]
So yes, they are sufficient, and the way GBP enters is instructive. The fair strike needs only the USD-collateralised EUR curve and SOFR:
it is the ordinary USD-CSA forward, as Q2 said it had to be. The GBP collateral appears only in the factor
$\df{\mathrm{SONIA}}/\df{[\mathrm{GBP},\mathrm{USD}]}=P^{(\gbp)}_{\usd}/P^{(\usd)}_{\usd}$, the ratio of the GBP-collateralised to the
USD-collateralised USD discount factor, which rescales the mark-to-market $K^{*}-K$. Equivalently, the EUR/GBP collateral basis is
triangulated from the two USD bases. Three caveats I would attach to that answer. The chaining uses the deterministic-spread
assumption; if the bases are stochastic and correlated with FX, each $P^{(k)}_i$ carries a convexity term that the five curves do not
determine. The triangulated EUR/GBP basis is a no-arbitrage value; where EUR/GBP basis swaps trade at a different level, a
GBP-collateralised EUR curve bootstrapped from them should replace the triangulation, and the formula above becomes the check against it.
And if a desk can invest collateral at something other than the market cash rate, the $y$'s become desk-specific: that is the
cheapest-to-deliver collateral question, not a curve question.

\section*{Part II: swap products in the valuation library}
I completed the six \code{TODO} blocks following the supplied FRA example. In \code{linear\_products.py}, \code{ProductRFRSwap} builds a
fixed leg and a compounded-overnight leg with \code{InterestRateStream}; both legs run from the adjusted effective date to the termination
date, their accrual schedules follow the overnight index's fixing calendar with the swap's payment business-day convention, and both pay
with the swap's payment offset and conventions. The floating leg carries the swap's compounding method and spread, which
\code{ProductOvernightIndexCashflow} adds outside the daily compounding. \code{ProductOvernightIndexBasisSwap} builds two such overnight
legs through one small helper, each on its own index, calendar and accrual period, with the spread on leg 1 only. The supplied
index-only and spread-only views of leg 1 reuse the leg-1 schedule, so leg 1 values to exactly their sum. In \code{product\_factory.py}
the two builders read everything but the dates and the quote from the convention; the shared date rule treats \code{1Y} as spot-starting,
\code{3Mx1Y} as a 3M forward start running for 1Y, and accepts explicit dates in either slot; the quote is the fixed rate for the
swap and the leg-1 spread for the basis swap. In \code{valuation\_engine.py} the two \code{calculate\_value} methods are signed sums of
the supplied leg engines' values and of their cash.

\textbf{Numbers.} Value date 2026-01-02, notional $N=1,000,000$ USD, trades from 2026-04-02 to 2027-04-02, funding on SOFR plus the
IFR spread. A useful identity for checking the floating legs by hand: when every observation period of the index calendar lies inside
the coupon, the daily factors $D(s_j)/D(e_j)$ telescope, so a compounded coupon pays $N\bigl(D(T_s)/D(T_e)-1\bigr)$ on the projection curve
plus $N\alpha s$ for the spread, discounted with the funding curve at the payment date. For the annual SOFR coupon this gives 4.2730\%, as in the notebook's cashflow table.
\vspace{2pt}{\centering\small
\begin{tabular}{@{}p{5.1cm}p{7.3cm}rl@{}}\toprule
Product & Leg values (USD) & PV (USD) & Fair level \\\midrule
FRA $3\times6$, receive 4\% (the supplied example) & fixed 9,937.00, floating 9,849.91 & 87.09 & \\
RFR swap, receive 4\% vs SOFR, 1Y/1Y & fixed 38,534.54, floating 41,164.22 & -2,629.68 & par rate 4.2730\% \\
Basis swap, SOFR + 10 bp vs Fed Funds, 3M/3M & leg 1 42,163.00 (index 41,182.75, spread 980.25), leg 2 33,716.53 & 8,446.47 & fair spread -76.17 bp \\
\bottomrule\end{tabular}\par}
\vspace{4pt}{\centering\small
\begin{tabular}{@{}lllrrrrrr@{}}\toprule
Start & End & Pays & $\alpha$ & SOFR cmpd \% & FF cmpd \% & $D_{\mathrm{fund}}$ & Leg 1 PV & Leg 2 PV \\\midrule
2026-04-02 & 2026-07-02 & 2026-07-07 & 0.252778 & 3.9649 & 2.7713 & 0.982228 & 10,092.66 & 6,880.75 \\
2026-07-02 & 2026-10-02 & 2026-10-06 & 0.255556 & 3.9652 & 3.6664 & 0.972239 & 10,100.34 & 9,109.54 \\
2026-10-02 & 2027-01-04 & 2027-01-06 & 0.261111 & 3.9656 & 3.6668 & 0.962191 & 10,214.35 & 9,212.31 \\
2027-01-04 & 2027-04-02 & 2027-04-06 & 0.244444 & 4.9614 & 3.6656 & 0.950167 & 11,755.64 & 8,513.93 \\
\bottomrule\end{tabular}\par}
\vspace{3pt}
The swap is negative for the fixed receiver because 4\% is below the 4.27\% that SOFR compounds to; the basis swap is worth 8,446.47
because Fed Funds projects 1 to 1.3 points below SOFR on these synthetic curves, far more than the 10 bp spread. Every number in the notebook agrees with the hand calculation above, and the 41 tests in
\code{tests/} check the same things: leg construction and dates, the factory's rules, the engines against the hand numbers, par rate
and fair spread, arithmetic compounding, same-day cash from fixings and past payments, the reports and the registry.
\end{document}
